\section{Pre-Registered Hypotheses and Their Outcomes}
\label{sec:prereg}
% Written 2026-09-23 from docs/CHECKLIST.md §1 (the ledger as registered on 2026-08-05) and the
% dated log entries that resolved each row. Verdicts are quoted from those entries; a row with no
% deciding entry is marked NOT RESOLVED rather than given a plausible outcome.

This project registered a hypothesis ledger before its main grids ran, and the study that resulted
is not the study that was planned. We report the ledger in full, including the rows that were
refuted, the row that was retired as ill-posed, and the two rows that were never resolved at all.
Table~\ref{tab:prereg} is that accounting. Three points in it are load-bearing for how the rest of
this paper should be read.

\paragraph{The original discriminator was refuted early, and the paper is a consequence of that.}
\textbf{H1c} as registered --- that a positive raw gain on the held-out obfuscator demonstrates
semantic invariance --- was retired within a day of the pilot, because an adapter trained on
\emph{clean} code reached the held-out family at least as well as one trained on obfuscated code
($0.414$ versus $0.384$). The raw gain measures task acquisition, not robustness. Its replacement,
\textbf{H1c-rev}, asks the same question relative to that clean-code control, and was refuted for
the condition it was first tested on ($-3.0$ points, CI $[-10.5, +3.9]$). Every claim in this paper
is stated against a clean-code-tuned control for this reason, and no raw difference from the
untuned model is ever reported as invariance.

\paragraph{Modularity was registered as a positive hypothesis and came back negative.}
\textbf{H2b} and \textbf{H2d} predicted that per-type adapters plus a router would beat monolithic
tuning, and that merges would retain most of the per-type gains. Neither held in the form
registered. The decisive measurement is that the two effects offset: merging \emph{costs}
$-3.13$ points (CI $[-4.78, -1.40]$) while the specialists \emph{contribute} $+2.47$
(CI $[+1.32, +3.70]$), for a net wash. \textbf{H2a}'s premise was confirmed and its consequence was
not: obfuscation type is almost perfectly surface-detectable --- a hard router dispatches at
$1.000$ accuracy --- and routing on that signal still bought nothing on the held-out family. The
present paper's RQ2 is the descendant of that result, which is why it compares composition
\emph{strategies} rather than arguing for one.

\paragraph{Two rows were never resolved, and we do not claim them.}
\textbf{H3} --- that the anchoring shift predicts which transfers succeed --- required a regression
of shift against transfer across all system $\times$ condition cells. That regression was never
run. What exists is the mechanism it presupposed: re-anchoring is real and causal
(\textbf{H3-causal}, confirmed by an identifier-knockout intervention, paired difference $+0.104$
nats, CI $[+0.011, +0.209]$), and an independent symbolic dead-code-elimination pass agrees with
it. That supports the mechanism and does \emph{not} establish the predictive claim, which we
therefore do not make. \textbf{HA1}, on alignment to human difficulty orderings, is likewise
unresolved: only the untuned correlation was read ($\rho = +0.132$, CI $[-0.140, +0.381]$), and the
pre/post change it was about was never measured.

\begin{table}[t]
\centering\small\setlength{\tabcolsep}{4pt}
\caption{\textbf{The pre-registered ledger, as registered and as it ended.} Registered
2026-08-05 before the main grids; verdicts are quoted from the dated log entry that decided each
row. \textsc{n/r} marks a row that no experiment resolved --- reported rather than dropped, because
a ledger that only lists its successes is not a ledger.}
\label{tab:prereg}
\begin{tabular}{@{}llp{3.1cm}@{}}
\toprule
ID & outcome & deciding evidence \\
\midrule
\textbf{H1a-trainable} & \textsc{supported} & self-gain $+27.3$ [$+11.0, +43.2$] \\
\textbf{H1a} family structure & \textsc{supported} & within-family $+33.3$ vs cross $+15.2$ \\
\textbf{H1b} cross-language & \textsc{partial} & JavaScript panel run; the registered contrast was not computed \\
\textbf{H1c} raw & \textsc{retired} & confounded with task acquisition \\
\textbf{H1c-rev} control-relative & \textsc{refuted} & $-3.0$ [$-10.5, +3.9$] on the held-out family \\
\textbf{H2a} detectable $\Rightarrow$ routable & \textsc{split} & dispatch $1.000$; no gain follows \\
\textbf{H2b} routed $\ge$ monolithic & \textsc{refuted} & as registered, at 1.5\,B \\
\textbf{H2c} conditioning vs capacity & \textsc{inconclusive} & recovery $0.26$--$0.33$, inside the reserved band \\
\textbf{H2d} merges retain $\ge 0.7$ & \textsc{refuted} & $-3.13$ cost against $+2.47$ contribution \\
\textbf{H3} shift predicts transfer & \textsc{n/r} & the regression was never run \\
\textbf{H3-causal} & \textsc{supported} & knockout $+0.104$ [$+0.011, +0.209$] nats \\
\textbf{HA1} human alignment & \textsc{n/r} & only the untuned $\rho$ was read \\
\bottomrule
\end{tabular}
\end{table}

\paragraph{Standing correctness checks.} The ledger was accompanied by a fixed list of
silent-failure checks to be cleared before any result was trusted: that splits partition by
program and never by row; that a tuned model's outputs actually differ from the untuned one's on
the same items; that one prompt builder serves training, both evaluation paths and attention
extraction; that the loss mask is verified on a real batch; that format-failure and truncation
rates are reported rather than assumed; that catastrophic forgetting is measured on a general
coding benchmark; and that headline numbers use the subset of programs on which every condition
succeeded. Each is enforced in code or reported in this paper. The forgetting check is the one that
proved hardest to make trustworthy, and Section~\ref{sec:threats} records why.
